\section{Discussion}
\label{sec:discussion}

\subsection{What the instrument buys an organisation}

The protocol converts an unfalsifiable assurance (``our scenarios describe the business rules'') into
three checkable quantities: the share of covered domain mutants the specification layer kills on its
own ($\mscov(\bdd)$); the unique uplift it contributes beyond the repository's other tests ($\dk$);
and the split between kills only it detects ($\bddunique$) and kills it also shares with those
tests ($\bddtotal$). The last two are the ones most organisations have never computed, and they
answer the question a manager actually asks --- \emph{is the Gherkin layer earning its maintenance
cost, or are the unit tests carrying the suite while the scenarios provide documentation?} A strong
overall mutation score together with $\dk \approx 0$ and $\bddunique \approx 0$ is a common and
previously invisible configuration: the suite is strong, and the specification layer is decorative.
The mirror-image configuration is equally worth naming: a layer with a healthy absolute score whose
$\bddunique$ is low is not broken, it is \emph{redundant}, and its maintenance cost buys no
detection that the rest of the suite does not already provide.

The coverage-conditioning rule matters here for a reason that is managerial as much as statistical. A
team told that its mutation score is 45\,\% will reasonably conclude that its tests are weak. If in
fact 40\,\% of the mutants were never executed by the spec layer --- because the scenarios exercise a
thin slice of the domain --- then the correct conclusion is that the scenarios do not cover the
decision logic, which is a coverage problem with a different and much cheaper remedy than
``write better assertions''. Conflating the two produces the wrong intervention.

\subsection{The finding that dominates: the population is thin}

The most consequential result of this work is not a correlation. It is that the population on which
such a correlation could be estimated at scale does not readily exist in public repositories. Of
\Nscreened{} screened repositories, \Neligible{} passed a static eligibility test that requires only
features, steps, a build, a licence, an application shape and a BDD runner. Of those, \Ngreen{}
produced a green, genuinely executing Gherkin baseline.

Three mechanisms explain this, and they are distinguishable in the attrition log. First, topic-tagged
BDD repositories are overwhelmingly \emph{tooling}: reporting agents, step libraries, formatters,
generators, boilerplates and demonstration projects. These are legitimate software, but they have no
business-rule domain to mutate --- a reporting agent has no pricing policy. Second, a large share of
application-shaped repositories are \emph{UI-driven}: their scenarios drive a browser, so a green
baseline requires a browser, a running application and often seeded data, none of which a
reproducible headless harness can assume. Third, ordinary \emph{dependency rot} removes a further
share: pinned dependency trees from two or three years ago no longer install.

For the research literature this has a sharp implication. A claim of the form ``we studied $N$ public
BDD suites'' should be read sceptically unless the authors report their funnel, because the naive
population --- repositories carrying a BDD topic tag --- is mostly not composed of the thing being
studied. Our contribution here is to make that funnel visible and categorised, so that future work can
either widen the frame deliberately or state its own selection more honestly.

\subsection{For Design Science: a negative result about deployment conditions}

In DSR terms, the artefact survived its feasibility evaluation and its population evaluation returned
a constraint rather than a correlation. This is not a failure of the artefact; it is a finding about
the context in which the artefact would be deployed at scale. Two readings are available and we do not
choose between them on the evidence we have.

The first reading is that the artefact is best deployed \emph{inside} organisations, on suites whose
build, services and data the team controls. There, the preconditions of Section~2 of the protocol can
actually be met, and the differential kill becomes a routine metric rather than a research instrument.
The public-ecosystem study then serves as a warning: suites that look like specifications from the
outside may not be executable at all, let alone oracular.

The second reading is that the public BDD ecosystem is genuinely dominated by scaffolding rather than
by specification-bearing applications, and that claims about ``BDD in practice'' drawn from public
repositories are therefore systematically skewed towards the scaffolding. Zampetti et al.\ already
observed that many projects use BDD frameworks without strictly applying BDD
\citep{zampetti2020demystifying}; our attrition data is consistent with that observation and adds a
mechanical explanation for why large-scale effectiveness studies are hard to run on the same
population.

\subsection{Implications for practice}

\begin{enumerate}
\item \textbf{Measure the unique uplift, not the total.} A suite-level mutation score cannot tell you
      whether the specification layer is doing any work. Run the three modes on one mutant set and
      report $\dk$, the share of mutants that only the specification layer kills. Comparing the
      BDD-only score with the full-suite score, as an earlier version of this protocol did, measures
      the opposite thing: it credits the specification layer with kills the unit tests already made.
\item \textbf{Treat ``not covered'' and ``not checked'' as different defects.} They have different
      remedies and different owners. The protocol's mandatory coverage conditioning exists to keep them
      apart.
\item \textbf{Track skip tags as a first-class metric.} The disabled-test literature shows that
      disabling is common and frequently permanent \citep{Kim2021Disabled}. A rising skip ratio is
      cheap to measure, needs no tooling, and is the earliest available signal that the specification
      layer is being withdrawn from service.
\item \textbf{Do not assume a scenario checks what its text says.} The appearance/function asymmetry is
      the whole problem. A scenario that reads well and asserts nothing is indistinguishable from a
      good one in review; it is distinguishable by mutation.
\end{enumerate}

\subsection{Implications for research}

The obvious next study is the one this paper could not run: a corpus of BDD suites drawn from
organisations rather than from public topic tags, where builds, services and data are controllable,
measured longitudinally against business-rule change events. Our instrument is designed to be reused
there, and the attrition analysis tells that study what to budget for. A second, cheaper direction is
to treat the differential kill as an outcome in mining studies that already have build access, and to
ask whether it declines before, after, or independently of the skip ratio --- a question about
precedence that requires more than two snapshots.

Finally, we would put the terminological point on record. ``Specification drift'' is not an
established construct; we define it here as a specialisation of architectural drift
\citep{PerryWolf1992,Rosik2010ArchDrift} to the executable-specification layer. If the term is useful,
it should acquire its own empirical grounding; if it is not, the underlying phenomenon --- a
specification artefact retaining its appearance while losing its discriminating power --- is still
real and still measurable, and the name matters less than the measurement.

\subsection{The organisational evaluation this study did not perform}

Design Science separates the artefact's \emph{technical} evaluation from its evaluation
\emph{in use}, and this paper delivers the first while only bounding the second. We measured the
deployment conditions at population scale, and we validated the indicators against known ground truth
on a suite we control; what we did not do is place the protocol in an organisation and observe a
decision being taken with it. That omission matters more for an Information Systems audience than for a
software-engineering one, because the IS contribution of such an instrument is not that it computes but
that it changes what a team attends to.

We therefore state the boundary rather than blur it. The claim this paper supports is that the
instrument is \emph{feasible}, that its indicators respond to injected drift (Section~\ref{sec:results}),
and that the public population cannot supply the evidence a single organisational case could. The claim
it does not support is that applying the protocol improves a team's specification practice. A
practitioner walkthrough --- one team, one repository, the protocol applied under observation, the
resulting decision recorded --- would test that second claim, and it is the natural next study rather
than a limitation we can argue away.
